% Generated draft: requires mathematical and source review.
\documentclass[11pt]{article}
\usepackage[margin=25mm]{geometry}
\usepackage{fontspec}
\setmainfont{Latin Modern Roman}
\usepackage{amsmath,amssymb,mathtools,unicode-math}
\setmathfont{Latin Modern Math}
\usepackage{enumitem,needspace,xurl,hyperref}
\setlength{\parindent}{0pt}
\setlength{\parskip}{5pt}
\setlength{\emergencystretch}{3em}
\newcommand{\N}{\mathbb{N}}
\newcommand{\Z}{\mathbb{Z}}
\newcommand{\Q}{\mathbb{Q}}
\newcommand{\R}{\mathbb{R}}
\newcommand{\C}{\mathbb{C}}
\newcommand{\F}{\mathbb{F}}
\newcommand{\A}{\mathbb{A}}
\newcommand{\PP}{\mathbb{P}}
\newcommand{\E}{\mathbb{E}}
\newcommand{\eps}{\varepsilon}
\newcommand{\dd}{\,\mathrm{d}}
\begin{document}
\section*{UnsolvedMath ID: 30005431}
\textbf{Source problem code:} catalogue entry
\subsection{Definitions and mathematical statement}
A \emph{skew brace} is a triple \((A,+,\circ)\) such that \((A,+)\) and \((A,\circ)\) are groups and
\[
 a\circ(b+c)=(a\circ b)+(-a)+(a\circ c)
 \qquad\text{for all }a,b,c\in A,
\]
where \(-a\) denotes the inverse of \(a\) in \((A,+)\). It is \emph{of abelian type} if \((A,+)\) is abelian. Its multiplicative group is \((A,\circ)\).

A finite group \(G\) is \emph{realizable} if there exists a finite skew brace \((A,+,\circ)\) of abelian type with \((A,\circ)\cong G\). Thus realizability concerns group isomorphism, and necessarily \(|A|=|G|\).

For a group \(G\), define its derived series by \(G^{(0)}=G\) and
\[
G^{(n+1)}=[G^{(n)},G^{(n)}],\qquad n\geq 0,
\]
where \([H,H]\) is the subgroup generated by all \(x^{-1}y^{-1}xy\) with \(x,y\in H\). The group \(G\) is \emph{solvable} if \(G^{(r)}=\{1\}\) for some integer \(r\geq 0\).

Determine a finite solvable group \(G\) such that \(G\) is not realizable and
\[
\text{for every finite solvable group }H,\qquad
|H|<|G|\ \Longrightarrow\ H\text{ is realizable}.
\]
In particular, identifying a candidate requires establishing both its failure of realizability and the realizability of every smaller finite solvable group. The target is one group of minimum order, rather than a classification of all groups attaining that order. The original excerpt asks for a ``minimal counterexample''; minimum order is the interpretation explicitly supplied by the cleaned statement.
\subsection{Short English statement}
Find a finite solvable group of the smallest possible order that cannot occur, up to isomorphism, as the multiplicative group of a skew brace whose additive group is abelian. Minimality requires every smaller finite solvable group to admit such a realization.
\subsection{Sources}
\begin{itemize}
\item Local UnsolvedMath registry: \texttt{data/unsolvedmath/problems.json}, permanent ID \(30005431\), exact source problem code \texttt{catalogue entry}; supplied complete source record, titled ``Smallest Solvable Group Excluded from Abelian-Type Skew Braces.'' The record gives proposed year \(2023\) and does not identify an individual proposer.
\item \emph{Mini-Workshop: Skew Braces and the Yang--Baxter Equation} (2023), Oberwolfach Reports, pp.~537--563, DOI \texttt{10.4171/owr/2023/9}, Problem~1 in the supplied excerpt. Original report: \url{https://doi.org/10.4171/owr/2023/9}.
\item Bachiller, \emph{Counterexample to a conjecture about braces}: \url{https://doi.org/10.1016/j.jalgebra.2016.01.011}. The supplied curation attributes counterexamples to the universal realization claim to this work.
\item Vendramin, described in the supplied record as a current skew-brace survey: \url{https://cris.vub.be/ws/portalfiles/portal/119985424/birkjour.pdf}. No more precise bibliographic description is supplied.
\item The original excerpt attributes to Etingof, Schedler and Soloviev, under its reference label [2], the result that the multiplicative group of a finite skew brace of abelian type is solvable. Its negative answer to the converse cites reference [1]. The excerpt does not include the numbered bibliography, so those reference labels cannot be independently matched here to full entries.
\item Literature assessment dated exactly \texttt{2026-08-22}: the supplied curation labels the minimum-counterexample problem open, reports that it remains posed in survey literature, and states that no minimal example was verified. This is supplied dated curation, not an independently verified literature assessment. The record also labels the statement \texttt{corrected\_verified} and reports direct checking of the official report; those are supplied provenance claims, not verification performed for this draft.
\end{itemize}
\end{document}
